% !TeX root = ../papers/nc-p3.tex
\section{整な余因子曲線の全特殊値における反対平方根の一意性}
\label{sec:integral-conductor-identifiability}
\begin{theorem}\label{thm:integral-conductor-identifiability}
被約有限窓の位置 $X$ に非恒等な全階数平方根 $N$ があり、行列式八次が整で、
余因子の正次元軌跡が整な一般被約曲線 $C$ であるとする。
このとき $X$ の非恒等な標識付き射影平方根はただ一つである。
別の根の行列式の整性、割線六次の滑らかさ、三重八次線形系の次元を仮定しない。
\end{theorem}
\begin{proof}[導手と二十次の場合]
既に $\deg C=20,21$ である。$D=V(\det N)$ の正規化を $\nu:S\to D$、入力束を $H$ とする。
二十次では $S=Q$、$\operatorname{coker}N=\nu_*\mathcal O_Q(1)$。
二十一次のBezout表示では $S$ は $T(sx)(z)=T(tx)(z)=0$ の母線束で、
$\operatorname{coker}N=\nu_*\pi^*\mathcal O(3)$ である。実際ある $z$ で二行が従属なら
像に平面が入り、各 $x$ の二四次式は独立なので射は有限、一般階数三から双有理である。
核座標による包含とEuler数の一致は定理\ref{thm:intrinsic-secant-scroll-root}の計算どおりである。

両場合に余因子イデアルは $D$ の導手と一致する。有限半局所な正規化環 $B$ 上で可逆層を自明化すると
元の環 $A$ 上の余核加群は $B$ となる。$\operatorname{Hom}_A(B,A)$ は分数体内の導手で、
その $B$ における全ての値の和も導手である。行列分解 $(N,\operatorname{adj}N)$ からこの値イデアルは
余因子の生成するイデアルである（$A$ 上で $N$ の行関係は随伴行列の行で生成される）。
有限双対により導手は $\nu_*(\omega_S\otimes\nu^*\omega_D^{-1})$、
従って正規化上の可逆イデアルで、$D$ 上で深さ二である。
商は純一次元Cohen--Macaulayで一般被約だから、導手スキームは $C$ そのものである。
特に孤立余成分はない。導手因子を $B_C$ とすれば $B_C\sim4H-K_S$。

二十個の全ての印で $I_C$ は極大イデアルの平方に入る
（階数一の対称行列の三次小行列式を展開する）。曲線局所環の埋込み次元は三で、
一般一次式で切った長さは少なくとも三である。よって各印の重複度は三以上。
$C$ の正規化上で印の極大イデアルの固定因子の次数はこの重複度である。
一般線形結合が各枝で最小付値をとり、その切断長は枝の位数の和になるためである。
任意の別の非恒等根 $N'$ の八次行列式は非零で、各印で三重、六次余因子は二重に消える。
$8\deg C<180$ より行列式は $C$ を含む。
二十次なら余因子を正規化へ制限して固定部を除いた次数は $120-120=0$ 以下である。
もし $C$ 上で一般階数三なら余因子が表す核点は一定で、
固定ベクトルとの $N'$ の積が $C$ 上で消える。十成分は独立なので非零二次式が二十印を通り矛盾する。
\end{proof}

\begin{proof}[二十一次の特異な割線六次にも残る線形系]
長さ二の二項因子 $q$ の平面に、二つの四次式 $f_x=T(sx),g_x=T(tx)$ を $q$ で割った余りをとる
四行四列の束写像を置く。標的は $E\oplus E$、
$0\to S_2\otimes\mathcal O(-1)\to S_4\otimes\mathcal O\to E\to0$ である。
行列式は六次であり、階数二以下の $q$ はない。そのような $q$ があれば入力の直線の全点で
二つの四次式に二次共通因子があるので、Bezout行列の階数二軌跡がその直線を含むからである。
従って零スキーム $P$ は $(q,x)$ の入射に同型である。射 $P\to C$ は有限で、
各 $x$ の繊維は $\gcd(f_x,g_x)$ の長さ二部分因子のスキームである。
一般の $C$ で最大公約数はちょうど二次である。これより行列式は恒等零ではない。
また $P$ は整である。実際入射上で $f_x=q f',g_x=q g'$ としたときBezout行列は
$q$ 倍の単射を $U_q:S_1\to S_3$ として
$U_q\operatorname{Bez}(f',g')U_q^t$ である。従って射は余因子スキーム $C$ に分解する。
階数二の一般点ではその核二平面、あるいは像二平面 $qS_1$ から二次式 $q$ がスキームとして復元される。
像が $qS_1$ を与える写像 $[q]\mapsto qS_1$ は閉埋込みである
（Plücker座標は $(sq)\wedge(tq)$、二点行列式で割れば $q\otimes q$）。
一般被約性は次のようにも直接見える。$C$ の関数体への上記入射の繊維は
固定二次因子の長さ二部分因子全体であり、一つの被約点である。$C$ 自体が被約なので
$P$ は一般に被約で、有限射により各曲線成分が $C$ を支配し、平面超曲面の純性から整となる。

$P$ 上で $H$ を入力超平面、$h$ を平面直線類とする。その次数は $21,6$。
残余 $f',g'$ は $H-h$ に値をとる二項二次式であり、結果式は $4(H-h)$ の次数六十の有効因子 $Z$ を与える。
各印で最大公約数は三次で、その長さ二部分因子スキームは長さ三である。
実際三次の中で二次因子を選ぶことは残る一次因子を選ぶことと同じで、
因子積写像の繊維として重根も含めて成立する。従って印の逆像イデアルの商長の合計は六十である。
$Z$ の局所式は上の二次Bezout行列の行列式であり、$U_q$ は局所的に左逆を持つので
元の行列の二次小行列式の生成イデアルが $(Z)$ のイデアルに一致する。
これらは印上で消える。包含と商長の等号により $I_X\mathcal O_P=\mathcal O_P(-Z)$ である。

$N'$ の六次余因子を $P$ に制限し、$2Z$ を除くと、次数 $126-120=6$ の\emph{可逆}層の切断になる。
この段階で $P$ を正規化しないことが重要である。
整な平面六次上で次数六の可逆層が三独立切断を持てば、それは $\mathcal O_P(h)$ である。
実際一切断の零因子 $W$ は平面の長さ六部分スキームで、随伴 $\omega_P=\mathcal O_P(3)$ と
Riemann--Rochから $h^0(\mathcal O_P(W))=1+h^1(\mathcal I_{W/\mathbb P^2}(3))$。
最大の直線交長を $r$ とすると、$3\le r\le5$ のとき直線残余の長さは三以下で二次式が値を分離し、
この $h^1$ は $\max(0,r-4)\le1$ 以下である。$r\le2$ のとき二点分を含む直線で二回残余を取れば、
最後の長さ二以下は一次式が分離するので同じ $h^1$ は零。従って $r=6$ で、
$W$ は直線切断に一致する。特に切断空間はちょうど三次元である。

もし $N'$ が $C$ 上で一般階数三なら、これらの切断による有理写像は核点の二次Veronese像に入る。
非定数ならその線形包は少なくとも平面である（Veronese像は直線を含まない）。
すると全切断系が平面六次そのものを写す。しかしVeronese像の定義二次式はこの六次上で消えれば
その平面上で恒等零となり、平面がVeronese像に入ってしまう。定数の場合は二十次と同じ二次式の矛盾である。
従って両次数で $N'$ の一般階数は $C$ に沿って二以下である。
\end{proof}

\begin{proof}[八次の強制と正規化加群の一意性]
$C$ の一般点で二行を消去するか行列式を展開すれば $\det N'$ は $I_C^2$ に入る。
もし $D$ 上で恒等零でなければ、$S$ でこの切断の零因子は $2B_C$ 以上である
（導手イデアルは $S$ 上で可逆、有限射は因子を点に潰さない）。
ところが $H^2=8$, $H B_C=40,42$ である。
二次曲面で $H=2L,K_S=-2L$、二十一次の母線束で $K_S=-2H+6F$, $HF=1$ を用いた。
$8H^2<2HB_C$ なので矛盾し、二つの八次式は比例する。

一般導手点の一次元局所環を $A$、正規化を $B$ と書く。
$A$ はGorenstein、導手が極大イデアルであるから $B/A$ の長さは一である
（有限双対で $B/A$ と $A/\operatorname{Hom}_A(B,A)$ の長さが一致する）。
$M'=\operatorname{coker}N'$ は階数一捩れなしで、この点で自由でない。
分数イデアルとして、その $B$ 倍を半局所Dedekind環の主イデアル性により $B$ に合わせる。
無限剰余体上で有限個の極大イデアルを同時に避ける生成元の一次結合があり、さらにこれで割れば
$A\subseteq M'\subseteq B$。長さ一と非自由性より $M'=B$ となる。
他の余次元一点でも $M'$ は正規化の加群である。深さ二によりこの作用は全体に延びる
（有限個の点の補集合から各局所切断を延長する）。従って $M'=\nu_*G$ で、
$G$ は $S$ 上の階数一反射層である。深さ二は有限射の各繊維上での局所コホモロジーの直和からも従う。

対称二次表示を $\mathcal O(-4)$ へ双対化すると
$\mathcal Ext^1(M',\mathcal O(-4))=M'(-2)$。有限双対と分数体上での線形性により
$2[G]=K_S+2H$ が $\mathrm{Cl}(S)$ で成り立つ。
滑らかな二次曲面と母線束のPicard群は自由、二次錐の類群も母線で生成される無限巡回群である
（$\mathbb F_2$ のPicard群を縮約切断で割る）。従って $G$ は先の正規化線束と同じである。
その自己準同型は定数だけで、極小対称表示の一意性
（定理\ref{thm:quadric-kernel-projection}の基底変更の証明）から $N'$ は出力合同とスカラーを除き $N$ に一致する。
\end{proof}
\path{computations/verify_integral_conductor_identifiability.py} は二次因子を持つBezout行列の多項式恒等式と
割線・導手・平面残余の次数評価を検算する。
